\section{问题二模型的建立与求解}

\subsection{问题二的分析}

问题二要求对附件中全部 50000 个 AI 推理与训练任务，在保持网络时延、截止时间、GPU 容量、IT 功率、设施功率和电网购售电边界等约束的前提下，通过跨区域迁移与开工时段调整，实现运行成本降低、碳排放减少与新能源利用率提升。

本问的核心挑战在于决策变量规模巨大：三类任务（AITraining、BatchInference、RealTimeInference）的完整合法时空候选约 $2.33\times10^8$ 个，传统精确求解方法难以直接处理。本文从任务结构、候选评价、求解收敛与多起点稳健性四个层面展开建模与求解。任务--准备--模型--算法--结果--解释六个环节的安排为：任务为带网络时延与截止时间约束的时空调度；准备为任务柔性差异、能源增量平衡与硬约束审计；模型为基于当前残余能源状态的动态边际 0--1 整数规划；算法为外层重复完整扫描直到收敛的 SFETA 构造式启发式；结果为 Cost-primary 与 Carbon-primary 双端点对照；解释为多起点路径依赖与极端等待的服务代价。

图\ref{fig:q2-flow} 给出本问的建模与求解总体流程。

\begin{figure}[!htbp]
\centering
\resizebox{\FigureFlowWidth}{!}{%
\begin{tikzpicture}[
  node distance=1.45cm and 1.35cm,
  box/.style={draw,rounded corners=2pt,minimum width=3.1cm,minimum height=1.0cm,align=center,fill=CumcmLightBlue},
  line/.style={-{Latex[length=2.5mm]},thick}
]
\node[box] (data) {任务与能源数据\\分析任务柔性差异};
\node[box,right=of data] (probe) {数据探针\\|R_i|、slack、GPU-hour};
\node[box,right=of probe] (model) {模型建立\\增量能量平衡\\碳感知时空调度};
\node[box,right=of model] (solve) {SFETA 求解\\重复完整扫描\\+ 收敛门禁};
\node[box,right=of solve] (result) {结果分析\\Cost/Carbon 双端点\\多起点稳健性};
\draw[line] (data)--(probe);
\draw[line] (probe)--(model);
\draw[line] (model)--(solve);
\draw[line] (solve)--(result);
\end{tikzpicture}%
}
\caption{问题二建模求解总体流程。}
\label{fig:q2-flow}
\end{figure}

\subsection{任务时空柔性分析}

\subsubsection{任务类型与基本参数}

设第 $i$ 个任务的到达时刻、来源区域、类型、GPU 需求数、持续时长、最晚完成时刻、最大允许网络时延分别为
\begin{equation}
a_i,\; o_i,\; k_i,\; g_i,\; p_i,\; d_i,\; L_i,
\label{eq:q2-task-params}
\end{equation}
其中 $p_i = \text{EstimatedDuration}_{\min}/60$（h），$k_i \in \{\text{RT},\text{Batch},\text{Training}\}$。

\subsubsection{时空可行域}

任务 $i$ 的空间可行域由网络时延约束决定：
\begin{equation}
R_i = \{\, r \mid \text{NetworkLatency}[o_i, r] \le L_i \,\},
\label{eq:q2-spatial-domain}
\end{equation}
时间可行域为
\begin{equation}
T_i = \begin{cases}
\{a_i\}, & k_i = \text{RT},\\[2pt]
\{\, s \mid a_i \le s,\; s + p_i \le \min(d_i, 2406) \,\}, & k_i \in \{\text{Batch},\text{Training}\}.
\end{cases}
\label{eq:q2-temporal-domain}
\end{equation}
完整合法时空候选集为 $C_i = R_i \times T_i$。RT 的合法域允许在全部 SLA 合法区域之间迁移，仅开工时刻固定为到达时刻。

\subsubsection{三类任务的柔性差异}

三类任务的时空柔性存在显著差异，图\ref{fig:q2-flexibility} 展示了各任务类型在空间可行区域数 $|R_i|$、时间余量 slack 和 GPU-hour 三个维度的分布。其中 slack 定义为
\begin{equation}
\text{Slack}_i = \min(d_i, 2406) - a_i - p_i.
\label{eq:q2-slack}
\end{equation}

\begin{figure}[!htbp]
\centering
\begin{subfigure}[t]{\FigureSingleWidth}
\centering
\caption{空间可行区域数 $|R_i|$ 分布}
\label{fig:q2-flex-region}
\end{subfigure}
\hspace{0.02\textwidth}
\begin{subfigure}[t]{\FigureSingleWidth}
\centering
\caption{GPU-hour 分布}
\label{fig:q2-flex-gpu}
\end{subfigure}
\caption{三类任务的时空柔性与资源重量对比。左：空间可行区域数分布；右：GPU-hour 分布。}
\label{fig:q2-flexibility}
\end{figure}

数据统计表明，RealTimeInference 的时间几乎刚性（$s_i = a_i$），空间可行域受 20 ms 时延限制，平均仅 2--3 个合法区域，且单任务 GPU-hour 最小。BatchInference 具有 5--6 个合法区域，slack 中位数约 1207 h，中等资源重量。AITraining 空间域全部开放，slack 中位数约 1184 h，且单任务 GPU-hour 中位数约 194.35，贡献了约 80\% 的总 GPU-hour 和约 87\% 的任务 AI 能耗。

这一差异构成了调度优先规则的数据基础：空间越受限、时间余量越小、资源重量越大的任务应优先安排，以保护最难调度的任务获得最佳资源位置。

\FloatBarrier

\subsection{能源层分析}

\subsubsection{AvailableRenewable 与基准弃电}

六区域的原始可用新能源功率 $\text{AvailableRenewable}[r,t]$ 逐时完全相同，因此区域间的新能源空间差异只能由基准运行状态下的弃电功率体现：
\begin{equation}
\text{Curtailment}_0[r,t] = \text{AvailableRenewable}[r,t] - \text{UsedRenewable}_0[r,t] - \text{RenewableCharge}_0[r,t] - \text{GridSell}_0[r,t].
\label{eq:q2-curtailment0}
\end{equation}
基准弃电在六区域间存在显著差异：每小时六区弃电最大值与最小值之差的中位数约 443.30 MW，最大约 580.52 MW。弃电最高的区域为 A 和 B，最低的区域为 E 和 F。

\subsubsection{电价与碳强度的关系}

各区域电价与碳强度的 Pearson 相关系数约 0.578，Spearman 相关系数约 0.613，且逐时 Pareto 分析表明 RegionE 在购电边际信号上始终占优。这意味着在弃电裕量耗尽、进入新增购电状态后，成本和碳目标在区域选择上总体同向，不会形成强烈的空间冲突。

\subsubsection{基准状态增量能量平衡}

为严格隔离 Q2 与 Q3/Q4 的决策自由度，Q2 保持基准状态的储能充放电与购售电策略不变，仅通过任务调度改变设施负荷。定义基准供负荷购电为
\begin{equation}
G_0[r,t] = \text{GridPurchase}_0[r,t] - \text{GridCharge}_0[r,t] \ge 0,
\label{eq:q2-g0}
\end{equation}
设施负荷增量为
\begin{equation}
\Delta L[r,t] = L[r,t](\mathbf x) - L_0[r,t],
\label{eq:q2-delta-l}
\end{equation}
其中 $L[r,t](\mathbf x) = \text{PUE}[r] \cdot \bigl(\text{NonAI\_IT\_Load}[r,t] + P_{\text{AI}}[r,t](\mathbf x)\bigr)$。

增量负荷 $\Delta L[r,t]$ 与基准弃电、新增购电的分段关系为
\begin{equation}
\Delta G[r,t] = 
\begin{cases}
0, & 0 \le \Delta L[r,t] \le \text{Curtailment}_0[r,t],\\
\Delta L[r,t] - \text{Curtailment}_0[r,t], & \Delta L[r,t] > \text{Curtailment}_0[r,t],\\
-\min\bigl(-\Delta L[r,t],\; G_0[r,t]\bigr), & \Delta L[r,t] < 0.
\end{cases}
\label{eq:q2-delta-g}
\end{equation}
该分段规则的含义是：新增负荷不超过基准弃电裕量时，全部由原弃电吸收，不增加购电和碳排放；超出部分才触发新增购电；负荷下降则优先减少原有供负荷购电。

更新后的弃电为
\begin{equation}
\Delta\text{Curtailment}[r,t] = \Delta G[r,t] - \Delta L[r,t],
\label{eq:q2-delta-curtailment}
\end{equation}
电网购电为 $\text{GridPurchase}[r,t] = \text{GridPurchase}_0[r,t] + \Delta G[r,t]$。

\subsection{模型建立}

\subsubsection{决策变量}

采用时间索引建模形式，定义二元决策变量
\begin{equation}
x_{i,r,s} = \begin{cases}
1, & \text{任务 } i \text{ 在区域 } r \text{、时刻 } s \text{ 开工},\\
0, & \text{否则}.
\end{cases}
\label{eq:q2-decision}
\end{equation}

\subsubsection{AI 算力负荷}

任务 $i$ 在区域 $r$、时刻 $t$ 的 AI IT 功率贡献为
\begin{equation}
P_{\text{AI}}[r,t](\mathbf x) = \frac{1}{\Delta t} \sum_{i,s} \alpha[k_i] \cdot g_i \cdot \omega(i,s,t) \cdot x_{i,r,s},
\label{eq:q2-ai-power}
\end{equation}
其中 $\alpha[k]$ 为任务类型 $k$ 的单位 GPU IT 功率（MW/GPU），$\Delta t = 1$ h，$\omega(i,s,t)$ 为小时重叠时长
\begin{equation}
\omega(i,s,t) = \max\bigl(0,\; \min(t+1,\; s+p_i) - \max(t,\; s)\bigr).
\label{eq:q2-overlap}
\end{equation}

\subsubsection{约束条件}

\textbf{唯一指派约束}：每个任务恰好选择一个 $(r,s)$ 组合：
\begin{equation}
\sum_{r\in R_i}\sum_{s\in T_i} x_{i,r,s} = 1,\quad \forall i.
\label{eq:q2-unique}
\end{equation}

\textbf{GPU-hour 容量约束}：
\begin{equation}
\sum_{i,s} g_i \cdot \omega(i,s,t) \cdot x_{i,r,s} \le \text{AvailableGPU}[r] \cdot \Delta t,\quad \forall r,t.
\label{eq:q2-gpu}
\end{equation}

\textbf{IT 功率约束}：
\begin{equation}
P_{\text{IT}}[r,t](\mathbf x) \le \text{MaxITPower}[r],\quad \forall r,t.
\label{eq:q2-it}
\end{equation}

\textbf{设施功率约束}：
\begin{equation}
L[r,t](\mathbf x) \le \text{MaxFacilityPower}[r],\quad \forall r,t.
\label{eq:q2-facility}
\end{equation}

\textbf{电网购售电边界}：
\begin{equation}
0 \le \text{GridPurchase}[r,t] \le \text{MaxGridImport}[r],\quad
\text{GridSell}_0[r,t] \le \text{MaxGridExport}[r],\quad \forall r,t.
\label{eq:q2-grid}
\end{equation}

\textbf{网络时延约束}：
\begin{equation}
\sum_{s} x_{i,r,s} \cdot \text{NetworkLatency}[o_i, r] \le L_i,\quad \forall i.
\label{eq:q2-latency}
\end{equation}

\textbf{截止时间约束}：任务 $i$ 必须在最晚完成时刻前结束
\begin{equation}
s_i + p_i \le \min(d_i, 2406),\quad \forall i.
\label{eq:q2-deadline}
\end{equation}

\textbf{非抢占与整时段约束}：任务在执行过程中区域固定、不可拆分；RT 必须即时开工但允许跨区域迁移
\begin{equation}
s_i = a_i,\quad k_i = \text{RT};\qquad
\text{执行过程中区域固定、不可拆分},\quad \forall i.
\label{eq:q2-nonpreempt}
\end{equation}

\subsubsection{目标函数}

运行成本
\begin{equation}
\text{Cost} = \sum_{r,t} \bigl[\text{GridPurchase}[r,t] \cdot \text{Price}[r,t] - \text{GridSell}_0[r,t] \cdot \text{SellPrice}[r,t]\bigr] \cdot \Delta t,
\label{eq:q2-cost}
\end{equation}
碳排放
\begin{equation}
\text{Carbon} = \sum_{r,t} \text{GridPurchase}[r,t] \cdot \text{CarbonIntensity}[r,t] \cdot \Delta t,
\label{eq:q2-carbon}
\end{equation}
新能源利用率
\begin{equation}
\eta_R = \frac{\sum_{r,t} \bigl(\text{UsedRenewable}[r,t] + \text{RenewableCharge}_0[r,t] + \text{GridSell}_0[r,t]\bigr) \cdot \Delta t}{\sum_{r,t} \text{AvailableRenewable}[r,t] \cdot \Delta t}.
\label{eq:q2-eta-r}
\end{equation}

Cost 与 Carbon 在候选评价中具有近似同向性（见数据探针分析），但二者并不完全退化。因此 Q2 采用 Cost-primary 为主目标，Carbon-primary 作为对照端点，Carbon 作为平行评价指标；CarbonIntensity 真实进入调度信号；不使用人工 Cost/Carbon 权重和。

\subsection{模型汇总}

综合以上分析，Q2 的完整数学模型可汇总为
\begin{equation}
\left\{
\begin{aligned}
&\text{决策变量：}\quad x_{i,r,s} \in \{0,1\},\; \forall i, r\in R_i, s\in T_i,\\
&\text{主目标：}\quad \min\; \text{Cost},\\
&\text{评价指标：}\quad \text{Carbon},\; \eta_R,\\
&\text{约束：}\quad \text{唯一指派、GPU-hour、IT 功率、设施功率、}\\
&\qquad\qquad \text{购售电边界、网络时延、截止时间、非抢占、RT 即时开工},\\
&\text{能源规则：}\quad \text{基准状态增量能量平衡},\\
&\text{边界：}\quad \text{固定基准储能充放电与购售电策略},\\
&\text{时域：}\quad \text{调度敏感时域 } [0, 2405];\; 2406 \text{ 仅做能源终端结算},\\
&\text{求解：}\quad \text{外层重复完整扫描至 } (J_{k-1}-J_k)/|J_{k-1}| < 10^{-3}.
\end{aligned}
\right.
\label{eq:q2-summary}
\end{equation}

\subsection{模型求解：SFETA 迭代收敛构造式启发式算法}

\subsubsection{算法设计思想}

面对约 $2.33\times10^8$ 的完整合法候选空间，本文设计 SFETA（Spatio-temporal Flexibility and Energy-aware Task Assignment）迭代收敛构造式启发式算法。该算法吸收 REPTA 的非迭代任务分配思想与 CarbonScaler 的边际资源分配思想，但针对本题的 GPU 规模固定、任务不可拆分、资源容量约束等特征进行特化。

独立原始数据复核发现，单遍全扫描存在明显路径依赖：第一遍完成后，前序任务的位移会改变 GPU/IT 与能源余量，从而可能为此前已处理任务重新创造更优 placement。因此 SFETA 在单遍构造式分配的基础上，外层增加重复完整扫描直到主目标相对改善率低于收敛容差，并将 last-pass fallback 计数与全部硬约束审计作为求解终止的联合门禁。算法的伪代码见算法\ref{alg:q2-sfeta}。

\begin{algorithm}[H]
\SetAlgoLined
\setstretch{0.92}
\SetAlFnt{\normalsize}
\SetArgSty{textnormal}
\SetKwInput{KwIn}{Input}
\SetKwInput{KwInit}{Initialize}
\SetKwInput{KwOut}{Output}
\SetKwFor{For}{for}{do}{end for}
\SetKwFor{ForEach}{for each}{do}{end for}
\SetKwIF{If}{ElseIf}{Else}{if}{then}{else if}{else}{end if}
\SetKw{KwRet}{return}
\SetKw{KwContinue}{continue}
\caption{SFETA：完整合法域动态边际任务分配（迭代收敛版）}
\label{alg:q2-sfeta}
\KwIn{任务集 $\mathcal I$，完整能源基准状态，网络/资源边界，目标模式 $m\in\{\mathrm{Cost},\mathrm{Carbon}\}$，相对收敛容差 $\epsilon$，最大完整扫描次数 $K_{\max}$}
\KwInit{由附件 source-local immediate 状态构造 $\mathbf x^{(0)}$；若该状态存在 GPU 超限，则先在完整合法域内执行可行性修复；$J^{(0)}\leftarrow+\infty$}
\BlankLine
\ForEach{任务 $i\in\mathcal I$}{
    构造空间合法域 $R_i=\{r:\mathrm{Latency}[o_i,r]\le L_i\}$\;
    构造时间合法域 $T_i$：RT 仅含 $a_i$；Batch/Training 使用全部满足 Arrival、LatestFinish 与 2406 的整数开工时刻\;
    记 $\Omega_i=R_i\times T_i$；若 $\Omega_i=\varnothing$ 则返回 FAIL\;
}
按 $|R_i|$ 升序 $\to$ Slack 升序 $\to$ GPU-hour 降序固定任务扫描顺序\;
\BlankLine
\For{$k\leftarrow1$ \KwTo $K_{\max}$}{
    从上一遍排程 $\mathbf x^{(k-1)}$ 重构 GPU/IT/Facility 与逐时能源状态\;
    $\mathrm{changed}\leftarrow0$，$\mathrm{fallback}\leftarrow0$\;
    \ForEach{排序后的任务 $i\in\mathcal I$}{
        读取任务 $i$ 的当前 placement $(r_i^{old},s_i^{old})$\;
        \textbf{暂时移除任务 $i$}，得到不含该任务的残余 workload 状态与能源状态 $S^{-i}$\;
        \ForEach{候选 $(r,s)\in\Omega_i$}{
            将任务 $i$ 按 fractional overlap 加回 $S^{-i}$，形成候选状态 $S^{-i}\oplus(i,r,s)$\;
            检查 GPU-hour、IT、Facility、GridImport、SLA、LatestFinish、finish$\le2406$；RT 额外检查 $s=a_i$\;
            若移除旧任务使某些能源小时暂时触及非负边界，则仅保留能够覆盖并修复这些 removal-dependent 小时的候选\;
            对通过硬约束的候选，计算动态边际响应
            \[
            \Delta E_{i,r,s}=\mathrm{EnergyResponse}(S^{-i}\oplus(i,r,s))-\mathrm{EnergyResponse}(S^{-i}),
            \]
            并由其得到 $\Delta\mathrm{GridPurchase}$、$\Delta\mathrm{Curtailment}$、$\Delta\mathrm{Cost}$ 与 $\Delta\mathrm{Carbon}$\;
        }
        得到全部合法候选集合 $\mathcal F_i$\;
        \If{$\mathcal F_i=\varnothing$}{
            对旧 placement 重新执行完整硬约束检查；若仍不可行则返回 FAIL，否则保留旧 placement 并令 $\mathrm{fallback}\leftarrow\mathrm{fallback}+1$\;
        }
        \Else{
            \If{$m=\mathrm{Cost}$}{按 $\Delta\mathrm{Cost}$ 升序、$\Delta\mathrm{Carbon}$ 次序选择候选\;}
            \Else{按 $\Delta\mathrm{Carbon}$ 升序、$\Delta\mathrm{Cost}$ 次序选择候选\;}
            对数值等价候选执行 tie-break：更早开工 $\to$ 保持 SourceRegion $\to$ 更低网络时延\;
            将最优候选重新插入当前状态；若 placement 改变则 $\mathrm{changed}\leftarrow\mathrm{changed}+1$\;
        }
    }
    由完整排程重新计算真实目标 $J^{(k)}$、硬约束审计和能源平衡残差\;
    \If{$k\ge2$ 且 $(J^{(k-1)}-J^{(k)})/|J^{(k-1)}|<\epsilon$ 且 $\mathrm{fallback}=0$ 且硬约束全部 PASS}{
        \textbf{break}\;
    }
}
\BlankLine
计算 $\mathrm{Cost}$、$\mathrm{Carbon}$、$\eta_R$、迁移/等待/网络时延及收敛轨迹；对 Cost-primary 可从不同初始可行排程重复运行用于路径依赖审计\;
\KwRet 调度方案 $\mathbf x$、主目标、平行评价指标和硬约束审计\;
\KwOut{各任务执行区域与开工时刻、Cost/Carbon/新能源利用率、迁移率、等待分布、收敛与初始化稳健性结果}
\end{algorithm}

\subsubsection{任务优先规则}

任务排序规则直接来自数据探针中三类任务的柔性差异，而非主观权重。排序键依次为
\begin{enumerate}
\item 空间可行区域数 $|R_i|$ 小者优先（空间柔性最小的任务最难迁移）；
\item slack $\text{Slack}_i$ 小者优先（时间余量最小的任务最紧迫）；
\item GPU-hour $g_i \cdot p_i$ 大者优先（资源重量最大的任务对全局影响最大）。
\end{enumerate}

该规则体现"先保护最难移动的任务，再处理大体量柔性任务"的调度逻辑，RealTimeInference 因时间刚性天然排在首位，AITraining 虽柔性最大但因 GPU-hour 极大而获得较高优先级。

\subsubsection{动态边际候选评价}

对每个任务的全体合法候选 $(r,s)$，基于当前残余调度状态计算真实的边际能源后果
\begin{equation}
\begin{aligned}
&\Delta\text{Curtailment}_{i,r,s},\quad
\Delta G_{i,r,s},\quad
\Delta\text{Cost}_{i,r,s},\quad
\Delta\text{Carbon}_{i,r,s},\\
&\text{同时检查 GPU/IT/Facility/GridImport/SLA/deadline 约束}.
\end{aligned}
\label{eq:q2-marginal}
\end{equation}

每次处理任务 $i$ 时，先从当前排程暂时移除 $i$，获得 removal 后的当前残余负荷与能源状态，再枚举完整合法 $(r,s)$ 域并对每个候选计算 $\text{EnergyResponse}(\text{current} + \text{candidate}) - \text{EnergyResponse}(\text{current})$。该口径避免了前序任务已消耗的弃电裕量被重复计入评价，也避免 removal 后的能源非负边界被忽略。

\subsubsection{等价候选的 tie-break}

当多个候选的边际能源后果在数值容差内等价时，按以下确定性规则选代表解：
\begin{enumerate}
\item 更早开工时刻 $s$；
\item 保持来源区域 $r = o_i$；
\item 更低网络时延 $\text{NetworkLatency}[o_i, r]$。
\end{enumerate}
该规则只处理等价方案，不覆盖 Cost/Carbon 主目标，不引入主观权重，也不会截断合法域。

\subsection{结果与分析}

SFETA 迭代收敛算法对全部 50000 个任务成功生成可行调度方案，所有硬约束均通过审计门禁。表\ref{tab:q2-metrics} 汇总了 Cost-primary 端点与附件参考状态的对比，Cost 与 Carbon 的相对节约均以 Q2 调度敏感时域口径（Hour 0--2405）报告。

\begin{table}[!htbp]
\centering
\caption{Q2 Cost-primary canonical 端点与附件参考状态对比（Q2 调度敏感时域 Hour 0--2405）}
\label{tab:q2-metrics}
\begin{tabular}{lccc}
\toprule
\rowcolor{CumcmLightBlue}
指标 & 附件参考状态 & SFETA Cost-primary & 变化 \\
\midrule
运行成本 / CNY & 1,801,786,872.69 & 1,555,371,304.13 & $-246,415,568.56$（$-13.68\%$）\\
碳排放 / tCO\textsubscript{2} & 2,045,367.48 & 1,826,774.35 & $-218,593.12$（$-10.69\%$）\\
新能源利用率 & 32.85\% & 37.29\% & $+4.44$ 个百分点 \\
跨区域迁移任务数 & 0 & 37,448 & 迁移率 $74.90\%$ \\
RT 迁移任务数 & 0 & 11,685 & 即时开工但跨区域 \\
最大 GPU 利用率 & 135.58\% & 100.00\% & 容量超限已消除 \\
last-pass fallback & -- & 0 & 全部任务由合法候选承载 \\
\bottomrule
\end{tabular}
\end{table}

\subsubsection{收敛轨迹}

Cost-primary 从 reference-start 出发完成 3 遍完整扫描，主目标单调下降并通过相对改善率门禁：
\begin{equation}
\text{pass1} \to \text{pass2} \to \text{pass3},\quad
\frac{J_2 - J_3}{|J_2|} = 5.77\times10^{-4} < 10^{-3}.
\end{equation}
收敛后 last-pass fallback = 0，全部硬约束 PASS。

\subsubsection{调度行为分析}

50000 个任务全部成功调度。表\ref{tab:q2-wait} 按任务类型给出迁移与等待分布。

\begin{table}[!htbp]
\centering
\caption{Q2 Cost-primary 端点的任务类型迁移与等待分布}
\label{tab:q2-wait}
\begin{tabular}{lcccccc}
\toprule
\rowcolor{CumcmLightBlue}
任务类型 & 任务数 & 迁移数 & 迁移率 & 平均等待/h & P95 等待/h & 最大等待/h \\
\midrule
AITraining & 16,559 & 12,059 & $72.824\%$ & 28.47 & 113 & 2010 \\
BatchInference & 16,717 & 13,704 & $81.976\%$ & 33.03 & 168.2 & 2016 \\
RealTimeInference & 16,724 & 11,685 & $69.870\%$ & 0 & 0 & 0 \\
\midrule
合计 & 50,000 & 37,448 & $74.90\%$ & 20.48 & 76 & 2016 \\
\bottomrule
\end{tabular}
\end{table}

注：上述逐类型统计来自 canonical Cost-primary 收敛排程；合计行的均值与 P95 按全部 50,000 个任务统计。RealTimeInference 全部保持 $s_i = a_i$，但允许在全部 SLA 合法区域之间迁移；canonical 端点中 11,685 个 RT 任务发生跨区域迁移，全部即时开工、零等待。

平均网络时延 26.75 ms，95\% 分位 78 ms，最大 82 ms，全部满足 SLA 要求。

\subsubsection{Carbon-primary 对照端点}

为检验 Cost/Carbon 目标结构，本文同步运行 Carbon-primary 端点（reference-start，3-pass 收敛）。表\ref{tab:q2-cmp} 给出 Cost-primary 与 Carbon-primary 的端点对比。

\begin{table}[!htbp]
\centering
\caption{Cost-primary 与 Carbon-primary 端点对照（reference-start，3-pass 收敛）}
\label{tab:q2-cmp}
\begin{tabular}{lcc}
\toprule
\rowcolor{CumcmLightBlue}
指标 & Cost-primary & Carbon-primary \\
\midrule
主目标 & Cost = 1,555,371,304.13 CNY & Carbon = 1,815,531.61 tCO\textsubscript{2} \\
对应副目标 & Carbon = 1,826,774.35 tCO\textsubscript{2} & Cost = 1,558,153,356.53 CNY \\
相对副目标代价 & -- & $+0.18\%$ Cost \\
相对副目标代价 & $+0.62\%$ Carbon & -- \\
\bottomrule
\end{tabular}
\end{table}

Carbon-primary 相比 Cost-primary 额外多花 2,782,052.40 CNY，仅占 Cost-primary 总成本 $0.18\%$；Cost-primary 相比 Carbon-primary 额外多排 11,242.75 tCO\textsubscript{2}，仅占 Carbon-primary 碳排 $0.62\%$。两者高度同向但并不完全退化，因此正文采用 Cost-primary 代表方案 + Carbon-primary 端点对照的组织方式，不构造人工 Cost/Carbon 加权和。

\subsubsection{多起点路径依赖与稳健性}

SFETA 为面向约 $2.33\times10^8$ 个合法候选的大规模动态边际构造式启发式，存在一定路径依赖。本文采用 reference-start（附件 source-local immediate 可行修复起点）与 previous-start（上一次扫描终态起点）双起点比较，约束容差均取 $10^{-3}$。表\ref{tab:q2-multistart} 给出两个起点的端点对比。

\begin{table}[!htbp]
\centering
\caption{Q2 多起点 Cost-primary 端点对比}
\label{tab:q2-multistart}
\begin{tabular}{lcc}
\toprule
\rowcolor{CumcmLightBlue}
指标 & reference-start & previous-start \\
\midrule
Cost / CNY & 1,555,371,304.13 & 1,576,174,462.12 \\
Carbon / tCO\textsubscript{2} & 1,826,774.35 & 1,838,076.57 \\
$\eta_R$ / \% & 37.29 & 37.72 \\
迁移率 / \% & 74.90 & 63.22 \\
平均等待 / h & 20.48 & 23.83 \\
P95 等待 / h & 76 & 94 \\
最大等待 / h & 2016 & 2239 \\
\bottomrule
\end{tabular}
\end{table}

previous-start 端点 Cost 相对 reference-start 高约 $1.34\%$。因此 Q2 选取全部硬约束满足且主目标更优的 reference-start 端点作为代表解；本文不声称算法对初始化不敏感，也不声称 50000 任务调度为全局整数最优。

\subsubsection{极端等待解释}

canonical Cost-primary 端点的最大等待为 2016 h。对该任务人工审计：TaskID 13051，BatchInference，Arrival = 69，Start = 2085，Finish = 2090.0667，LatestFinish = 2406，AvailableSlack $\approx$ 2331.93 h，实际仅使用约 $86.45\%$ 合法 slack，Latency = 22 ms $<$ 80 ms。前若干个极端等待任务全部为 Batch/Training，无 RT；全部满足 SLA、LatestFinish 与 $\text{Finish}_i \le 2406$。

长等待是 Cost-primary 利用官方超长时间柔性的服务代价，反映了"能源主目标 vs. 即时性"的权衡；它不是时间索引或 deadline bug。Q2 不人为加入 24/48 h 最大等待窗来掩盖这一代价；Cost--QoS 的正式联合权衡留给 Q4 处理。

\subsubsection{Hour 2406 跨问结算口径}

当前 Q2 求解代码的经济核算为 Hour 0--2405；Q3 能源模型包含 Hour 0--2406，且 Hour 2406 只做能源终端结算、禁止计算任务占用。两侧的公共常数差异约为 $1.267\times10^5$ CNY。本节给出的绝对 Cost/Carbon/$\eta_R$ 数字均显式标注为 Q2 调度敏感时域口径（Hour 0--2405），不与 Q3 B0 直接并表比较；最终跨问绝对对比等待建模手完成统一 accounting/registry 后再冻结。

\subsubsection{柔性捕获率定义}

为进一步评估调度方案对附件理论可利用柔性的利用程度，定义柔性捕获率
\begin{equation}
\text{FlexCapture}_{\text{cost}} = \frac{\text{实际成本节约}}{\text{理论成本节约上界}},\quad
\text{FlexCapture}_{\text{carbon}} = \frac{\text{实际碳减排}}{\text{理论碳减排上界}},
\label{eq:q2-flex-capture}
\end{equation}
其中上界来自连续松弛下的 removal-side 分析。新 canonical 端点下 FlexCapture 的具体数值待统一 registry 完成后补入；本节只给出定义与计算口径，不引用旧单遍的百分比。

\subsection{模型检验}

\subsubsection{硬约束审计}

全部硬约束在 last-pass 均通过：
\begin{itemize}
\item GPU-hour、IT 功率、设施功率：违规 0；
\item 电网购售电上下界：违规 0；
\item 网络时延 SLA：违规 0；
\item 逐任务 $\text{Finish}_i \le \text{LatestFinishHour}_i$：违规 0；
\item 统一 $\text{Finish}_i \le 2406$ 且 Hour 2406 无计算占用：通过；
\item RT 强制 $\text{StartHour} = \text{ArrivalHour}$：违规 0；
\item last-pass fallback 计数：0；
\item 统一能量平衡最大残差约 $2.00\times10^{-4}$ MW $< 10^{-3}$ MW；
\item AvailableRenewable 分解最大残差约 $2.00\times10^{-4}$ MW；
\item baseline AI IT 重构最大误差约 $3.33\times10^{-7}$ MW；
\item overlap 守恒、全部能源变量非负：通过。
\end{itemize}

\subsubsection{Baseline 复现检验}

当 $\Delta L[r,t] \equiv 0$（即不改变任何任务的调度位置）时，增量能量平衡精确复现附件基准的能源结果，验证了增量核算口径的一致性。

\subsubsection{参数敏感性分析}

参数敏感性分析需对 Cost-primary canonical 端点重算后写入正文；旧 ±5\% 扰动结论来自旧单遍数值，不复用于当前结果。本节保留待重算占位，不给出具体百分比。

\subsubsection{柔性可用性上界验证}

连续松弛下的流体可行性分析表明，在放松任务级时间约束与不可拆分约束后，附件存在足够的时空弃电与容量承接全部 Batch 和 Training 工作量，且不触发任何新增购电。这说明 Q2 的调度结果主要受任务级 temporal/deadline/nonpreemptive 约束和构造式算法的限制，而非总量意义上的绿色承载不足。
